\section{Gaussian transport and a matrix surplus}\label{tr:section}

We first work with a probability density $e^{-V}$ on $\mathbb R^m$ satisfying
\begin{equation}\label{tr:smooth-assumptions}
 V\in C^\infty(\mathbb R^m),\qquad
 0<c_- I\le D^2V\le c_+I<\infty.
\end{equation}
The constants $c_-,c_+$ may depend on the density.  Write
$C=\operatorname{Cov}(\mu)$, where $d\mu=e^{-V}\,dy$, and let
$\gamma_m$ be standard Gaussian measure.  Throughout this section,
$\mathbb E$ denotes integration against $\gamma_m$.

There are two objectives.  We choose affine coordinates in which a
reparametrization of the transport Jacobian has mean zero.  We then use
convexity of $V$ to control derivatives of that reparametrization.  The
resulting estimate contains an additional nonnegative term measuring
separation between eigenvalues; retaining this term is essential in the
matrix argument that follows.

\subsection{The Gaussian parametrization}

Let $\Phi$ and $\varphi$ denote the distribution function and density of a
standard one-dimensional Gaussian.  Define, for $a\in\mathbb R$,
\begin{equation}\label{tr:scalar-functions}
 \begin{gathered}
 t(a)=-\log\Phi(-a),\qquad q(a)=t'(a)=\frac{\varphi(a)}{\Phi(-a)},
 \qquad d(a)=q(a)-a,\\
 k(a)=\frac{t(a)}{q(a)},\qquad b(a)=k'(a),\qquad l(a)=b'(a).
 \end{gathered}
\end{equation}
The map $t$ transports a standard Gaussian to the exponential law of
mean one: $\Phi(-G)$ is uniform on $(0,1)$ when $G$ is standard Gaussian.
Thus $q$ is the derivative of the scalar transport associated with the
entropy benchmark $h=1$, $\operatorname{Var}=1$.

The identities
\begin{equation}\label{tr:scalar-identities}
 q'=qd,\qquad b=1-kd,\qquad l=k-q+ab
\end{equation}
follow by differentiation.  In particular, $q-k=ab-l$.  We record the
elementary properties that permit us to use $q$ as a change of variable.

\begin{lemma}\label{tr:q-properties}
The function $q$ is a smooth increasing bijection from $\mathbb R$ to
$(0,\infty)$.  Moreover,
\[
 d>0,\qquad 0<q'<1,\qquad q''\ge0,
 \qquad 0<1-qd\le d^2.
\]
\end{lemma}
\begin{proof}
Let $X_a$ have the conditional law of $G-a$ given $G>a$, and write
$\mathbb E_a$ for expectation under this law.  Gaussian integration
by parts gives
\[
 \mathbb E_a X_a=d(a),\qquad
 \operatorname{Var}(X_a)=1-q(a)d(a).
\]
Since $X_a>0$ almost surely, $d>0$ and hence $q'=qd>0$.
Its strictly positive variance gives $q'<1$.

For completeness, this variance is at most the square of its mean.
The survival function
\[
 S_a(u)=\frac{\Phi(-a-u)}{\Phi(-a)},\qquad u\ge0,
\]
satisfies $S_a(u+v)\le S_a(u)S_a(v)$, since its hazard rate $q(a+u)$
is increasing.  Integrating over $u,v\ge0$ yields
\[
 \frac12\mathbb E_a X_a^2
 =\int_0^\infty uS_a(u)\,du
 \le\left(\int_0^\infty S_a(u)\,du\right)^2=d(a)^2.
\]
Consequently $1-qd=\operatorname{Var}(X_a)\le d^2$.  Differentiating
$q'=qd$ now gives
$q''=q(d^2+qd-1)\ge0$.  Finally, $q(a)\to0$ as $a\to-\infty$,
whereas $q(a)>a$ for $a>0$.  These limits prove surjectivity.
\end{proof}

We shall also need two positivity bounds and a quantitative growth estimate.
Their proof, including
the rational bounds on the compact interval and estimates on both tails,
is given in Appendix~\ref{cert:verification}.
Put
\begin{equation}\label{tr:M-definition}
 D_0(a)=\frac1{d(a)},\qquad M(a)=D_0(a)^2(1-b(a)^2).
\end{equation}
A lower bound for the growth of $M$ in $\log q$ will let us compare
uniform averages in $q$ with uniform averages in $\log q$ in the
surplus estimate.
Define a nondecreasing step function $m_0$ by the following table.  At a
finite endpoint we take the value in the interval to its right.
\begin{equation}\label{tr:bins}
\begin{array}{c|rrrrrr}
 &O&P_1&R_1&U_1&V_1&T_1\\\hline
 a&(-\infty,-3)&[-3,-2)&[-2,-1)&[-1,\tfrac12)&[\tfrac12,4)&[4,\infty)\\
 m_0(a)&0&.19&.38&.68&.68&.68
\end{array}
\end{equation}
All terminating decimals in this paper denote exact rational numbers.

\begin{lemma}[Scalar bounds]\label{tr:scalar-basic}
The functions in \eqref{tr:scalar-functions} satisfy, on $\mathbb R$,
\begin{equation}\label{tr:scalar-basic-equations}
 0<b<\frac12,\qquad l>0,\qquad
 \frac{d\log M}{d\log q}
 :=\frac{M'}{Md}\ge m_0.
\end{equation}
\end{lemma}

For a real symmetric matrix, scalar functions will always be interpreted
by spectral functional calculus.  If $f$ is differentiable, its divided
difference is
\begin{equation}\label{tr:divided-difference}
 f[u,v]=\begin{cases}(f(u)-f(v))/(u-v),&u\ne v,\\f'(u),&u=v.
 \end{cases}
\end{equation}
The classical Daleckii--Krein formula
\cite{DaleckiiKrein1965} states that, in an eigenbasis of a symmetric
matrix $A$, the derivative of $f(A)$ in a symmetric direction $Z$ has
entries $f[a_i,a_j]Z_{ij}$; see also
\cite[Section~2.4]{Noferini2017}.  This formula
will allow us to keep the off-diagonal entries of matrix derivatives.

\subsection{Transport and the choice of coordinates}

Brenier's theorem gives a convex potential $\phi$ such that
$(\nabla\phi)_\#\gamma_m=\mu$; its gradient is unique up to null sets
\cite{Brenier1991,McCann1995}.  For $m\ge2$ we use the full-space
regularity theorem of Cordero-Erausquin and Figalli
\cite[Theorem~1.1]{CorderoErausquinFigalli2019}, which extends
Caffarelli's interior theory \cite{Caffarelli1992} to the unbounded
domains considered here.  It applies with both domains equal to $\R^m$, because both densities
are smooth and locally bounded away from zero.  For $m=1$, the same
smoothness follows directly by expressing monotone transport through
the two distribution functions.  Thus $\phi$ is smooth in every dimension.
Caffarelli's contraction theorem \cite{Caffarelli2000} and its inverse
bound then imply that the Jacobian
\begin{equation}\label{tr:Jacobian}
 J=D^2\phi
\end{equation}
satisfies
\begin{equation}\label{tr:Jacobian-bounds}
 c_+^{-1/2}I\le J\le c_-^{-1/2}I.
\end{equation}
For the lower bound, apply the contraction theorem to the inverse
transport.  We use its quantitative form: a source potential with
Hessian at most $aI$ and a target potential with Hessian at least $bI$
give a transport Lipschitz constant at most $\sqrt{a/b}$; see also
\cite[Theorem~1]{ChewiPooladian2023}.

Define the symmetric matrix field
\begin{equation}\label{tr:A-definition}
 A=q^{-1}(J).
\end{equation}
The bounds \eqref{tr:Jacobian-bounds} place the spectrum of $A$ in a
compact interval.  Our preferred coordinates are characterized by
$\mathbb EA=0$.  They need not be the usual isotropic coordinates.

\begin{proposition}[Affine normalization]\label{tr:normalization}
Let $\mu=e^{-V}dy$ satisfy \eqref{tr:smooth-assumptions}.  There exists a
positive definite symmetric matrix $P$ such that the Brenier map from
$\gamma_m$ to $P_\#\mu$ has Jacobian $J_P$ satisfying
\[
 \mathbb E q^{-1}(J_P)=0.
\]
The transformed density again satisfies
\eqref{tr:smooth-assumptions}, with possibly different positive constants.
\end{proposition}
\begin{proof}
We use a homotopy from the Gaussian and show that its zeros stay in a
compact subset of the positive definite cone.  For $0\le\theta\le1$,
define
\[
 W_\theta(y)=(1-\theta)|y|^2/2+\theta V(y),\qquad
 d\mu_\theta=Z_\theta^{-1}e^{-W_\theta(y)}\,dy,
\]
where $Z_\theta$ is the normalizing integral.
There are constants $a_-,a_+>0$, independent of $\theta$, with
$a_-I\le D^2W_\theta\le a_+I$.  These densities have uniformly
Gaussian tails.  Their covariance matrices $C_\theta$ depend
continuously on $\theta$, are positive definite, and therefore satisfy
\begin{equation}\label{tr:homotopy-covariance}
 c_0I\le C_\theta\le C_0I
\end{equation}
for fixed $c_0,C_0>0$.

For a positive definite symmetric $P$, let $J_{\theta,P}$ be the
Jacobian of the transport to $P_\#\mu_\theta$.  Define a map with
values in the vector space of symmetric matrices by
\[
 F(\theta,P)=\mathbb E q^{-1}(J_{\theta,P}).
\]
Contraction in both directions gives
\begin{equation}\label{tr:homotopy-spectral}
 \frac{\lambda_{\min}(P)}{\sqrt{a_+}}I
 \le J_{\theta,P}\le\frac{\|P\|_{\mathrm{op}}}{\sqrt{a_-}}I.
\end{equation}

We first verify continuity of $F$.  Restrict $P$ to any compact subset
of the positive definite cone.  The transports have common Lipschitz
constants by \eqref{tr:homotopy-spectral}; their values at zero are
bounded because their means and second moments are bounded.  Every
convergent sequence $(\theta_j,P_j)$ thus has locally uniformly
convergent subsequences of transport maps.  Their limits push forward
$\gamma_m$ to the limiting target and remain gradients of convex
functions, so Brenier uniqueness identifies the limit.  Write $J_j,J$
for the corresponding Jacobians.  They converge weak-* locally and
obey common positive spectral bounds.

The change-of-variables equation expresses $\log\det J_j$ in terms
of the target potential composed with its transport map.  It therefore
converges locally uniformly to $\log\det J$.  On a fixed compact
spectral interval, strict concavity of $\log\det$ gives a constant
$c>0$ such that
\[
 \log\det J_j\le\log\det J+
 \operatorname{tr}\bigl(J^{-1}(J_j-J)\bigr)-c|J_j-J|_{\mathrm F}^2.
\]
Integration on any ball, followed by weak-* convergence, proves
$J_j\to J$ in local $L^2$.  Functional calculus then gives local
$L^2$ convergence of $q^{-1}(J_j)$.  Their common spectral bounds
control the Gaussian tails, proving continuity of $F$.

We next bound all zeros, uniformly in $\theta$.  Suppose
$\mathbb EA=0$, where $A=q^{-1}(J_{\theta,P})$.  Convexity of $q$
and Jensen's inequality on the expected spectral measure of each unit
vector give
\[
 \mathbb E J_{\theta,P}\ge q(0)I.
\]
Gaussian integration by parts gives
$\mathbb E(x_r\partial_i\phi)=(\mathbb E J_{\theta,P})_{ir}$.
Orthogonal projection of the centered components of $\nabla\phi$
onto the span of $x_1,\ldots,x_m$ consequently yields
\[
 P C_\theta P\ge(\mathbb E J_{\theta,P})^2\ge q(0)^2I.
\]
Together with \eqref{tr:homotopy-covariance}, this bounds $P^{-1}$
uniformly.  The lower estimate in \eqref{tr:homotopy-spectral} now
gives a uniform lower bound $A\ge a_*I$, where we may take $a_*\le0$.
Since $q'\le1$,
\[
 q(a)\le q(a_*)+a-a_*\quad(a\ge a_*),
 \qquad
 \mathbb E\operatorname{tr}J_{\theta,P}
 \le m\bigl(q(a_*)-a_*\bigr).
\]
On the other hand, Gaussian Poincar\'e and
\eqref{tr:homotopy-spectral} imply
\[
 c_0\|P\|_{\mathrm{op}}^2
 \le\operatorname{tr}(P C_\theta P)
 \le\mathbb E\operatorname{tr}J_{\theta,P}^2
 \le\frac{\|P\|_{\mathrm{op}}}{\sqrt{a_-}}
       \mathbb E\operatorname{tr}J_{\theta,P}.
\]
This also bounds $P$ uniformly.

Choose a bounded open set in the space of symmetric matrices whose
closure lies in the positive definite cone and whose interior contains
every zero just bounded.  The map $F(\theta,\cdot)$ has no zeros on its
boundary.  At $\theta=0$, the target is $N(0,P^2)$, so
$J_{0,P}=P$ and $F(0,P)=q^{-1}(P)$.  Its unique zero is $q(0)I$;
its derivative there is multiplication by the positive scalar
$1/q'(0)$.  Its Brouwer degree is therefore nonzero.  Homotopy
invariance of degree gives a zero at $\theta=1$, as required.
\end{proof}

The quantity
$m+\tfrac12\log\det\operatorname{Cov}(\mu)-h(\mu)$ is unchanged
by an invertible affine transformation.  Indeed, a linear map $P$
adds $\log|\det P|$ to entropy and $2\log|\det P|$ to the covariance
log determinant.  We may therefore impose the normalization from
Proposition~\ref{tr:normalization}.  From now on $V,J,A,C$ refer to
this transformed density and transport, and
\begin{equation}\label{tr:mean-zero}
 \mathbb EA=0.
\end{equation}

\subsection{The transport equation and its energy}

We shall differentiate the Jacobian and integrate the resulting equation.
The twice-differentiated change-of-variables equation is part of
Caffarelli's transport method \cite[Section~5]{Caffarelli2000}.
The next lemma establishes the finite energy needed to do this without
additional decay assumptions on higher derivatives of $V$.
For a matrix field $F$, set
\[
 |F|_{\mathrm F}^2=\operatorname{tr}(F^{\mathsf T}F),\qquad
 \|F\|^2=\mathbb E|F|_{\mathrm F}^2,\qquad
 \|\partial F\|^2=\sum_{r=1}^m\|\partial_rF\|^2.
\]
The same convention sums over any displayed derivative index of a
collection of matrix fields.

\begin{lemma}[Jacobian equation and finite energy]\label{tr:energy}
Under \eqref{tr:smooth-assumptions}, put $J_r=\partial_rJ$ and
\[
 \mathcal L u=\operatorname{tr}(J^{-1}D^2u)
             -\nabla V(\nabla\phi)\cdot\nabla u.
\]
Then
\begin{equation}\label{tr:Jacobian-equation}
 \mathcal L J_{ij}
 =-\delta_{ij}+(JD^2V(\nabla\phi)J)_{ij}
   +\operatorname{tr}(J^{-1}J_iJ^{-1}J_j).
\end{equation}
Moreover $\|\partial J\|<\infty$.  The matrix fields $A$ and $k(A)$
belong to Gaussian $H^1$, and the integrations by parts below with
smooth matrix functions of $J$ are justified by compact cutoffs.
\end{lemma}
\begin{proof}
The density equation is
\begin{equation}\label{tr:density-equation}
 \log\det J=V(\nabla\phi)-|x|^2/2-(m/2)\log(2\pi).
\end{equation}
Differentiating twice gives \eqref{tr:Jacobian-equation}.  The divergence
of the cofactor matrix of a Hessian vanishes.  Applying this identity
to $J$, and then differentiating \eqref{tr:density-equation} once,
shows that
\[
 \mathcal L u=e^{|x|^2/2}\operatorname{div}
       \bigl(e^{-|x|^2/2}J^{-1}\nabla u\bigr).
\]
Thus, for smooth $u$ and compactly supported smooth $v$,
\begin{equation}\label{tr:divergence-form}
 \mathbb E(v\mathcal Lu)
 =-\mathbb E\langle\nabla v,J^{-1}\nabla u\rangle.
\end{equation}

Choose $0<\lambda\le\Lambda<\infty$ such that
$\lambda I\le J\le\Lambda I$.  Taking the trace in
\eqref{tr:Jacobian-equation}, and using $D^2V\ge0$, gives
\[
 \mathcal L(\operatorname{tr}J)
 \ge-m+\Lambda^{-2}\sum_r|J_r|_{\mathrm F}^2.
\]
For a compactly supported cutoff $0\le\chi\le1$, integration by parts
and $|\nabla\operatorname{tr}J|\le\sqrt m\,|\partial J|$ imply
\begin{align*}
 \Lambda^{-2}\mathbb E\chi^2|\partial J|^2
 &\le m+2\lambda^{-1}\sqrt m\,
             \mathbb E\bigl(\chi|\nabla\chi|\,|\partial J|\bigr)\\
 &\le m+\tfrac12\Lambda^{-2}\mathbb E\chi^2|\partial J|^2
       +2m\Lambda^2\lambda^{-2}\mathbb E|\nabla\chi|^2.
\end{align*}
Let $\chi$ tend to one through standard radial cutoffs with gradients
bounded by a constant divided by their radii.  Fatou's lemma proves
$\|\partial J\|<\infty$.

All eigenvalues of $J$ lie in $[\lambda,\Lambda]$.  The Fr\'echet
derivatives of $q^{-1}$ and of $k\circ q^{-1}$ are bounded on that
spectral interval, so $A,k(A)\in H^1(\gamma_m)$.  The same applies to
any smooth scalar function on this interval.  For such a matrix
multiplier, insert a cutoff into \eqref{tr:divergence-form} and let it
tend to one.  The terms containing its gradient vanish by
Cauchy--Schwarz and the energy bound; the remaining differentiated
terms are bounded by a constant times $1+|\partial J|^2$.
This justifies the asserted integrations by parts.
\end{proof}

\subsection{Keeping the eigenvalue separation}

Set $Z_r=\partial_r A$.  At each point diagonalize $A$, and write
$a_i$ for its eigenvalues and $\lambda_i=q(a_i)$ for those of $J$.
For real $a,b$, define
\begin{equation}\label{tr:surplus-weight}
 \begin{gathered}
 \delta(a,b)=|\log q(a)-\log q(b)|,\qquad
 m_*(a,b)=m_0(\min(a,b)),\\
 r(a,b)=\frac{1+m_*(a,b)}{16}
    \left(1+\frac{\delta(a,b)^2}{30}
             +\frac{\delta(a,b)^4}{1680}\right).
 \end{gathered}
\end{equation}
We use the abbreviations $\delta_{ij}=\delta(a_i,a_j)$ and
$r_{ij}=r(a_i,a_j)$.  The quantity
\begin{equation}\label{tr:surplus-definition}
 \mathcal C=\mathbb E\sum_{i,j,r}
           r_{ij}\delta_{ij}^2(Z_r)_{ij}^2
\end{equation}
is independent of the choice of orthonormal basis within repeated
eigenspaces.  In particular, the factor $\delta_{ij}^2$ vanishes when
$a_i=a_j$.

\begin{proposition}[Matrix surplus]\label{tr:surplus}
For the transport satisfying \eqref{tr:smooth-assumptions}, with the
notation above,
\begin{equation}\label{tr:surplus-inequality}
 \mathbb E\operatorname{tr}k(A)^2
 \ge\|\partial A\|^2-\|\partial k(A)\|^2+\mathcal C.
\end{equation}
\end{proposition}
\begin{proof}
The proof has two parts.  Symmetry of the third derivatives of $\phi$
first replaces the transport equation by averages of the scalar
function $M$.  Comparing arithmetic and logarithmic averages then
produces the term $\mathcal C$.

The identity $qk=t$ will match the entropy term in the covariance
completion of Section~\ref{mat:section}; we first estimate the derivative
cost of $k(A)$. Put $F(\lambda)=k(q^{-1}(\lambda))^2$. Multiply
\eqref{tr:Jacobian-equation} by $F(J)$ in trace and integrate by parts.
Lemma~\ref{tr:energy} justifies this operation.  At a fixed point
choose an orthonormal eigenbasis of $J$ and rotate all three indices
of the symmetric tensor
$\Theta_{ijr}=\partial_rJ_{ij}=\partial_i\partial_j\partial_r\phi$
into that basis.  The quadratic term on the right of the Jacobian
equation has trace contraction
\[
 \sum_i F(\lambda_i)\operatorname{tr}(J^{-1}J_iJ^{-1}J_i)
 =\sum_{i,j,r}\frac{F(\lambda_i)}{\lambda_j\lambda_r}\Theta_{ijr}^2,
\]
where we used the symmetry of $\Theta$.  The term obtained by integration
by parts is
\[
 \sum_r\lambda_r^{-1}\operatorname{tr}\bigl((\partial_rF(J))J_r\bigr)
 =\sum_{i,j,r}\frac{F[\lambda_i,\lambda_j]}{\lambda_r}\Theta_{ijr}^2.
\]
The remaining term is
$\operatorname{tr}(F(J)JD^2VJ)\ge0$.
Discarding it gives
\begin{equation}\label{tr:initial-multiplier}
 \mathbb E\operatorname{tr}k(A)^2
 \ge\mathbb E\sum_{i,j,r}\Theta_{ijr}^2
 \left(\frac{F(\lambda_i)}{\lambda_j\lambda_r}
             +\frac{F[\lambda_i,\lambda_j]}{\lambda_r}\right).
\end{equation}
Only pointwise changes of basis are made here; no eigenvectors are
differentiated.

Regard $M$ as a function of $\lambda=q(a)$.  The identities
\eqref{tr:scalar-identities} give
\begin{equation}\label{tr:lambdaF}
 \frac{d}{d\lambda}(\lambda F(\lambda))
 =k^2+\frac{2kb}{d}
 =\frac{1-b^2}{d^2}=M.
\end{equation}
For two positive numbers $u,v$, let $M^{\mathrm{ar}}_{uv}$ be the
uniform average of $M$ on the interval between $u$ and $v$, with its
continuous interpretation when $u=v$.  With the weight
$\Theta_{ijr}^2/(\lambda_i\lambda_j\lambda_r)$, symmetrization in
$i,j$ changes the coefficient in \eqref{tr:initial-multiplier} to
\[
 \frac{\lambda_i+\lambda_j}{2}\,
 M^{\mathrm{ar}}_{\lambda_i\lambda_j}.
\]
In the full sum we can replace this coefficient by
$\lambda_rM^{\mathrm{ar}}_{\lambda_i\lambda_j}$, decreasing its
value.  To see this, order three eigenvalues as $u\le v\le w$.
Lemma~\ref{tr:scalar-basic} makes $M$ increasing, so
\[
 M^{\mathrm{ar}}_{uv}\le M^{\mathrm{ar}}_{uw}
                         \le M^{\mathrm{ar}}_{vw}.
\]
The ordered coefficients $u+v-2w$, $u+w-2v$, and $v+w-2u$ sum to
zero.  Their scalar product with these three ordered averages is
nonnegative, by the rearrangement inequality.  Full symmetry of
$\Theta$ permits precisely this averaging over permutations of the
three indices.  We have proved
\begin{equation}\label{tr:arithmetic-energy}
 \mathbb E\operatorname{tr}k(A)^2
 \ge\mathbb E\sum_{i,j,r}
 \frac{\Theta_{ijr}^2}{\lambda_i\lambda_j\lambda_r}\,
 \lambda_r M^{\mathrm{ar}}_{\lambda_i\lambda_j}.
\end{equation}

We now fix a pair of eigenvalues.  Bars in the rest of the proof denote
uniform averages in $\log\lambda$ over their interval.  Set
\[
 x=\delta_{ij}/2,\qquad w_0=\frac{x}{\sinh x},\qquad m_*=m_*(a_i,a_j),
\]
with $w_0=1$ at $x=0$.  Since $da/d\log\lambda=D_0$,
the divided-difference formula shows that the coefficients of
$(Z_r)_{ij}^2$ and $(\partial_r k(A))_{ij}^2$, respectively, when
written with the same weight as in \eqref{tr:arithmetic-energy}, are
\begin{equation}\label{tr:derivative-coefficients}
 \lambda_r w_0^2\overline{D_0}^{\,2},\qquad
 \lambda_r w_0^2\overline{bD_0}^{\,2}.
\end{equation}
Both $D_0(1+b)$ and $D_0(1-b)=k$ are increasing: $d'=qd-1<0$,
and $b'=l>0$, $k'=b>0$.  The covariance inequality for two increasing
functions on an interval therefore yields
\begin{equation}\label{tr:logarithmic-chebyshev}
 \overline M\ge
 M_1:=\overline{D_0}^{\,2}-\overline{bD_0}^{\,2}
 \ge\tfrac34\overline{D_0}^{\,2}.
\end{equation}

We also need a quantitative comparison of the two averaging measures.
On the logarithmic interval, \eqref{tr:scalar-basic-equations} says
that $M(\lambda)\lambda^{-m_*}$ is increasing.  Exponential weighting
and the same covariance inequality give
\begin{equation}\label{tr:average-comparison}
 \frac{M^{\mathrm{ar}}_{\lambda_i\lambda_j}}{\overline M}
 \ge\frac{\operatorname{shc}((1+m_*)x)}
          {\operatorname{shc}(x)\operatorname{shc}(m_*x)},
 \qquad \operatorname{shc}(x)=\frac{\sinh x}{x}.
\end{equation}
Indeed, after centering the logarithmic interval, the left side is
$\mathbb E_0(e^sM)/[\mathbb E_0e^s\,\mathbb E_0M]$, where
$\mathbb E_0$ is its uniform average.  Under the measure proportional
to $e^{m_*s}$, the increasing factor $Me^{-m_*s}$ can only increase
the mean of $e^s$.  Taking $M=e^{m_*s}$ gives the displayed ratio.

Here is a convenient explicit lower bound for this ratio.  Define
$A_0(x)=x\coth x-1$, with $A_0(0)=0$.  The identity
\[
 x\coth x+\frac{x^2}{\sinh^2x}\ge2
\]
implies both $A_0(x)\ge1-w_0^2$ and that $A_0(x)/x^2$ is
nonincreasing for $x>0$.  To verify the identity, clear the denominator
and expand
$\tfrac{x}{2}\sinh(2x)+x^2-\cosh(2x)+1$ in powers of $x$;
the coefficients vanish through degree four and are nonnegative
thereafter.  Since $0\le m_*\le1$, we obtain
$A_0(m_*x)\ge m_*^2A_0(x)$.  The addition formula for $\sinh$ now gives
\begin{align}
 \frac{\operatorname{shc}((1+m_*)x)}
          {\operatorname{shc}(x)\operatorname{shc}(m_*x)}
 &=1+\frac{m_*A_0(x)+A_0(m_*x)}{1+m_*}\notag\\
 &\ge1+m_*(1-w_0^2).
 \label{tr:hyperbolic-comparison}
\end{align}
Every formula here extends continuously to $m_*=0$ or $x=0$.

Combining \eqref{tr:logarithmic-chebyshev}--\eqref{tr:hyperbolic-comparison},
the amount by which the coefficient in
\eqref{tr:arithmetic-energy} exceeds the difference of the two
coefficients in \eqref{tr:derivative-coefficients} is at least
\[
 \lambda_r(1+m_*)(1-w_0^2)M_1.
\]
Relative to the coefficient of $(Z_r)_{ij}^2$, this is at least
\[
 \frac34(1+m_*)(w_0^{-2}-1)
 \ge\frac{1+m_*}{16}
    \left(\delta_{ij}^2+\frac{\delta_{ij}^4}{30}
                            +\frac{\delta_{ij}^6}{1680}\right).
\]
The last inequality is the Taylor expansion of
$(\sinh x/x)^2$, whose omitted coefficients are positive.
Substituting in \eqref{tr:arithmetic-energy} and summing proves
\eqref{tr:surplus-inequality}.
\end{proof}

We have obtained two inputs for the entropy argument: the matrix field
$A$ has no constant Gaussian component, and convexity supplies
\eqref{tr:surplus-inequality}.  The next section combines these facts
with a Gaussian covariance identity and separates the first chaos of
$A$ from its higher chaoses.
